\section{Data Development}\label{sec_data}

\subsection{Data Acquisition}

The Los Angeles Planning Department's website maintains robust public documentation of City Planning Commission meetings. For each meeting, the agenda, minutes, and any supplemental documents relevant to the meeting (letters from the public, traffic assessments, architectural reports, etc.) are available for download as PDF files.\footnote{As of June 18th, 2026, these documents are available at the URL: \url{https://planning.lacity.gov/about/commissions-boards-hearings}.} We downloaded these documents for all City Planning Commission meetings from May 10th, 2018 to February 2nd, 2026. This resulted in documentary data for \gn{NumberOfMeetings} meetings, covering \gn{NumberOfAgendaItems} agenda items, with \gn{NumberOfSupplementalDocs} supplemental documents, spanning \gn{PageCount} pages of PDF documents. Download occurred twice: once on April 10th, 2025, and a second time on May 15th, 2026.

Since we are primarily interested in the regulatory approvals process, we limit our attention to the agenda items that require a decision from the board. These are identified by agenda items titled according to their Planning Department case numbers, which have a standardized format of ``[CASE PREFIX]-[YEAR]-[SERIAL NUMBER]-[CASE SUFFIXES]''. The prefix identifies the decision-making body initially responsible for approval, and the case suffixes identify requested entitlements and other case features. Other agenda items include items like ``Director's Report'' and ``General Public Comment'' which typically do not require any decision-making on the part of the board. Altogether, there were \gn{NumberOfCases} agenda items requiring a decision from the board, as identified by their Planning Department case numbers.

A typical agenda item is a request from a developer to approve a development plan that goes beyond what the site's zoning designation would allow. Figure \ref{fig_example_agenda_item} shows an example of an agenda item. The case number is DIR-2019-6048-TOC-SPR-WDI-1A. This was a project that was initially approved by the Director of Planning (DIR). The suffixes indicate that the project is requesting bonuses under the Los Angeles Transit Oriented Communities program (TOC), the project requires a discretionary Site Plan Review (SPR), and the project was granted a waiver of dedications and improvements (WDI). However, this previously approved plan was appealed (1A), and both the appeal and the project proposal are now being considered by the City Planning Commission. In addition to the information contained in the case number, the published agenda contains other data such as the Council District that the project is located in and other specific details about the project proposal.

Figure \ref{fig_example_minutes_item} shows the associated minutes for this agenda item. From the minutes, we can see that the appeal was granted in part and denied in part. The CPC upheld the Director of Planning's previous decision, but additional conditions were applied, thus allowing the project to move forward as long as the developer adheres to the new conditions. This outcome---the partial granting of an appeal and the approval of a project with conditions or modifications---is common but not the only kind of outcome. Sometimes, the applicant's requested actions are granted in their entirety without additional conditions or modifications. Rarely, the requested actions are denied entirely. A more common occurrence than denial is that the CPC puts the decision off to a later date. 
%We will discuss the distribution of motion outcomes and voting patterns later in Section \ref{sec_descriptive_statistics}.

Figures \ref{fig_example_support_letter} and \ref{fig_example_oppose_letter} show examples of letters submitted by the public in support of and in opposition to the above project. The support letter emphasizes how the project will ease traffic, reduce air pollution, and increase housing availability. The opposition letter emphasizes concerns about displacement and how the proposed units will be unaffordable to current residents of the neighborhood. These letters typify the kinds of concern expressed by community residents in this dataset; however, the number of letters that this project attracted is atypical (DIR-2019-6048 just happens to be a particularly controversial project). We will discuss the distribution of support and opposition letters across projects later in Section \ref{sec_methodology}.

\subsection{Data Extraction}

The documentary data provides a wealth of information about CPC cases and their outcomes. However, the information is locked within textual documents that are difficult to extract using traditional methods. For example, traditional NLP methods based on token and pattern matching would have a hard time comparing the agenda to the minutes and determining whether the requested actions were approved, partially approved, approved with conditions, denied, or whether the decision was postponed to a future date. 

To extract usable features from the textual content more robustly, we make use of Anthropic's \texttt{claude-opus-4-6} generative language model. The model can be asked to read the text of the agenda, read the text of the minutes, and to extract structured information about the decisions of the Commission. For example, it can be asked explain what the result of committee's proposed motion was in terms of its implications for the development project---whether the project was approved, approved with conditions, denied, or whether the decision was postponed to a later date. The methodology and prompts we used to perform the data extraction are described in detail in the appendix, and are robust to using alternative language models. The full data extraction pipeline, including all prompts and scraping code, is publicly available at \url{https://github.com/ed-kung/lacpc}. In this section, we will simply focus on the data features that were extracted from the published texts.

\paragraph{Agenda items.} For each agenda item, we extracted the following information from the published agenda: 
\begin{itemize}[noitemsep, topsep=0pt]
\item Agenda item number (used primarily for identification);
\item Item title (for cases requiring a decision, this is always a Planning Department case number);
\item Short LLM-generated summary of the agenda item's content;
\item The Council District(s) to which the item applies;
\item If the case is a development proposal, the physical dimensions of the proposal where applicable (number of units, square footage, height, parking spaces);
\item List of requested actions.
\end{itemize}

\paragraph{Minutes.} For each agenda item, we extract the following information from the published minutes text:
\begin{itemize}[noitemsep, topsep=0pt]
\item Short LLM-generated summary of the deliberations and the motion that was ultimately voted on;
\item Implication of the motion for the proposed project: whether the requested actions were approved, approved with minor conditions and/or modifications, approved with major conditions and/or modifications, denied, or whether deliberations were continued to a future date;
\item Result of the vote (whether the motion passed or failed);\footnote{Note that a motion passing is not the same as the project getting approved, nor is a motion failing the same as the project getting denied. A Member may move to deny the project's requested actions, or move to accept an appeal of a previously approved project, in which case the motion passing implies a denial of the project. These nuances highlight why LLMs are helpful in the data extraction process.}
\item Vote tallies: the number of ayes, nays, absences, abstentions, and recusals.
\end{itemize}

\paragraph{Supplemental documents.} For each supplemental document, we extract the following information from its text:
\begin{itemize}[noitemsep, topsep=0pt]
\item Type of document: whether it is a public comment expressing support or opposition to the case, or whether it was technical documentation attached to the project or an administrative or procedural document;
%\item Type of author: whether the author of the document is an individual, an advocacy group, a consultant, a lawyer, a developer, or a public official;
\item Which agenda item(s) it references;
\item Short LLM-generated summary of the document contents;
\item Support or opposition: For each agenda item that the document references, whether the document supports, opposes, or is neutral towards the development proposal in the referenced agenda item(s);
\item Reasons for support or opposition: For each referenced agenda item that the document supports or opposes, a list of mentioned reasons for the support or opposition.\footnote{The list of possible reasons that the LLM was asked to select from are shown in Table \ref{tab_reasons}.}
\end{itemize}

